\par\Needspace{7\baselineskip}\bigskip
\begin{center}
\pdfbookmark[2]{Gauss an Olbers, Göttingen, 31. März 1818}{pallas-brief-18180331}
\textsc{Gauss} an \textsc{Olbers}, Göttingen, 31. März 1818.
\end{center}
\noindent\hspace*{\parindent}. . . . .\quad Da meine Tafeln [der Störungen], welche Pallas abseiten des Jupiter erleidet, jetzt ganz vollendet und die Elemente allen 12 bisherigen Oppositionen möglichst genau angepasst sind, so habe ich Hrn. \textsc{Encke} ersucht, sämmtliche Beobachtungen von 1802 nach diesen Elementen zu berechnen. Obgleich Sie den Planeten erst etwas nach der Opposition entdeckten, so wird sich jetzt, da die Elemente so scharf bestimmt sind, jene doch vielleicht noch aus den Beobachtungen ableiten lassen.\quad . . . . .
\par
